\subsection{The eta series and the signed hexagonal Gaussian}
\label{sec:eta-signed-mellin}
Let $\chi_{12}(n)=(12/n)$ be the primitive even character of conductor
$12$, with values $1,-1,-1,1$ at $1,5,7,11$ modulo $12$.
Euler's pentagonal identity gives the exact unary expansion
\begin{equation}\label{eq:eta-unary-chi12}
 \eta(\tau)=\sum_{n\geq1}\chi_{12}(n)q^{n^2/24}.
\end{equation}
Indeed its pentagonal exponent is
$1/24+k(3k-1)/2=(6k-1)^2/24$, and taking the absolute values of
$6k-1$, for $k\in\mathbb Z$, enumerates the positive integers prime to
$6$ exactly once with sign $\chi_{12}$.
Thus, with the Mellin Dirichlet convention that a term $a_\nu q^\nu$
contributes $a_\nu\nu^{-s}$,
\begin{equation}\label{eq:eta-mellin-dirichlet}
 L(\eta,s)=24^sL(2s,\chi_{12}),\qquad
 \int_0^\infty\eta(iy)y^{s-1}\,dy
   =\Gamma(s)\left(\frac{12}{\pi}\right)^s L(2s,\chi_{12}).
\end{equation}
Termwise integration is initially valid for $\Re s>1/2$.
The transformation of $\eta$ shows exponential decay, up to a power,
at both integration endpoints, so the completed integral extends
entirely.  If the sum in \eqref{eq:eta-unary-chi12} were instead taken over
all $n\in\mathbb Z$, it would need a factor $1/2$.
Alp\"oge's theorem places the nontrivial zeros of $L(\eta,s)$
on $\Re s=1/4$.

For the signed Gaussian, use the Euclidean hexagonal lattice
$L=\mathbb Z+\omega\mathbb Z$ with
$\omega=(-1+\sqrt{-3})/2$ and $|a+b\omega|^2=a^2-ab+b^2$.
There are four cosets of $L$ in $\frac12L$.  Define
\[
 \epsilon(v)=\begin{cases}-1,&v\in L,\\+1,&v\in\frac12L\setminus L,
 \end{cases}
 \qquad
 G_\epsilon(t)=\sum_{v\in\frac12L}\epsilon(v)e^{-\pi t|v|^2}.
\]
The convention fixes the lattice scale: the quadratic form is $|v|^2$,
without division by $\Im\omega$.  Writing
$\Theta_L(t)=\sum_{v\in L}e^{-\pi t|v|^2}$ gives
\[
 G_\epsilon(t)=\Theta_L(t/4)-2\Theta_L(t),\qquad G_\epsilon(\infty)=-1.
\]
For $\Re s>1$, absolute convergence and termwise integration yield
\begin{align}
 \frac{\pi^s}{\Gamma(s)}
   \int_0^\infty\bigl(G_\epsilon(t)+1\bigr)t^{s-1}\,dt
 &=\sum_{0\ne v\in\frac12L}\epsilon(v)|v|^{-2s}\notag\\
 &=(4^s-2)Z_L(s)\notag\\
 &=6(4^s-2)\zeta(s)L(s,\chi_{-3}).
 \label{eq:signed-hexagonal-mellin}
\end{align}
The last equality follows by identifying nonzero Eisenstein integers
modulo their six units with nonzero ideals; the Dedekind zeta function of
$\mathbb Q(\sqrt{-3})$ is $\zeta(s)L(s,\chi_{-3})$.
Both sides then continue meromorphically.
The extra elementary factor has zeros
\[
 s=\frac12+\frac{\pi i k}{\log2},\qquad k\in\mathbb Z.
\]
If the area-normalized quadratic form $|v|^2/\Im\omega$ is used instead,
the right side acquires the nonvanishing factor $(\Im\omega)^s$.
